\documentclass[10pt,a4paper]{amsart}
\usepackage{amsmath,amssymb,amsthm,graphicx}
% Retained minimal preamble for the Stacks Project chapter more-morphisms.tex (More on Morphisms).
% Verbatim excerpt of the chapter preamble: Xy-pic setup (native preamble lines 14-18) and the native
% theorem declarations with the mathematical macros (native preamble lines 154-217). Copied mechanically;
% nothing added, removed or redefined. The wrapper class and the amsmath/amssymb/amsthm/graphicx base
% packages are supplied by the assembler, so the class/externaldocument part of the native preamble is omitted.
% For commutative diagrams we use Xy-pic
\usepackage[all]{xy}

% We use 2cell for 2-commutative diagrams.
\xyoption{2cell}
% Theorem environments.
%
\theoremstyle{plain}
\newtheorem{theorem}[subsection]{Theorem}
\newtheorem{proposition}[subsection]{Proposition}
\newtheorem{lemma}[subsection]{Lemma}

\theoremstyle{definition}
\newtheorem{definition}[subsection]{Definition}
\newtheorem{example}[subsection]{Example}
\newtheorem{exercise}[subsection]{Exercise}
\newtheorem{situation}[subsection]{Situation}

\theoremstyle{remark}
\newtheorem{remark}[subsection]{Remark}
\newtheorem{remarks}[subsection]{Remarks}

\numberwithin{equation}{subsection}

% Macros
%
\def\lim{\mathop{\mathrm{lim}}\nolimits}
\def\colim{\mathop{\mathrm{colim}}\nolimits}
\def\Spec{\mathop{\mathrm{Spec}}}
\def\Hom{\mathop{\mathrm{Hom}}\nolimits}
\def\Ext{\mathop{\mathrm{Ext}}\nolimits}
\def\SheafHom{\mathop{\mathcal{H}\!\mathit{om}}\nolimits}
\def\SheafExt{\mathop{\mathcal{E}\!\mathit{xt}}\nolimits}
\def\Sch{\mathit{Sch}}
\def\Mor{\mathop{\mathrm{Mor}}\nolimits}
\def\Ob{\mathop{\mathrm{Ob}}\nolimits}
\def\Sh{\mathop{\mathit{Sh}}\nolimits}
\def\NL{\mathop{N\!L}\nolimits}
\def\CH{\mathop{\mathrm{CH}}\nolimits}
\def\proetale{{pro\text{-}\acute{e}tale}}
\def\etale{{\acute{e}tale}}
\def\QCoh{\mathit{QCoh}}
\def\Ker{\mathop{\mathrm{Ker}}}
\def\Im{\mathop{\mathrm{Im}}}
\def\Coker{\mathop{\mathrm{Coker}}}
\def\Coim{\mathop{\mathrm{Coim}}}

% Boxtimes
%
\DeclareMathSymbol{\boxtimes}{\mathbin}{AMSa}{"02}

%
% Macros for moduli stacks/spaces
%
\def\QCohstack{\mathcal{QC}\!\mathit{oh}}
\def\Cohstack{\mathcal{C}\!\mathit{oh}}
\def\Spacesstack{\mathcal{S}\!\mathit{paces}}
\def\Quotfunctor{\mathrm{Quot}}
\def\Hilbfunctor{\mathrm{Hilb}}
\def\Curvesstack{\mathcal{C}\!\mathit{urves}}
\def\Polarizedstack{\mathcal{P}\!\mathit{olarized}}
\def\Complexesstack{\mathcal{C}\!\mathit{omplexes}}
% \Pic is the operator that assigns to X its picard group, usage \Pic(X)
% \Picardstack_{X/B} denotes the Picard stack of X over B
% \Picardfunctor_{X/B} denotes the Picard functor of X over B
\def\Pic{\mathop{\mathrm{Pic}}\nolimits}
\def\Picardstack{\mathcal{P}\!\mathit{ic}}
\def\Picardfunctor{\mathrm{Pic}}
\def\Deformationcategory{\mathcal{D}\!\mathit{ef}}
\newenvironment{reference}{}{}
\newenvironment{slogan}{}{}
\newenvironment{history}{}{}
\usepackage{cleveref}
\crefname{originallemma}{original lemma}{original lemmas}
\Crefname{originallemma}{Original lemma}{Original lemmas}
\makeatletter
\newlabel{lemma-flat-locus-base-change}{{lemma-flat-locus-base-change}{source}{}{sourcepointer-1}{}}
\newlabel{lemma-flat-locus-base-change@cref}{{[originallemma][0][]lemma-flat-locus-base-change}{[1][0][]source}{}{}{}}
\makeatother
\crefname{originallemma}{original lemma}{original lemmas}
\Crefname{originallemma}{Original lemma}{Original lemmas}
\makeatletter
\newlabel{morphisms-lemma-finite-presentation-permanence}{{morphisms-lemma-finite-presentation-permanence}{source}{}{sourcepointer-2}{}}
\newlabel{morphisms-lemma-finite-presentation-permanence@cref}{{[originallemma][0][]morphisms-lemma-finite-presentation-permanence}{[1][0][]source}{}{}{}}
\makeatother
\crefname{originallemma}{original lemma}{original lemmas}
\Crefname{originallemma}{Original lemma}{Original lemmas}
\makeatletter
\newlabel{theorem-criterion-flatness-fibre}{{theorem-criterion-flatness-fibre}{source}{}{sourcepointer-3}{}}
\newlabel{theorem-criterion-flatness-fibre@cref}{{[originallemma][0][]theorem-criterion-flatness-fibre}{[1][0][]source}{}{}{}}
\makeatother
\crefname{originallemma}{original lemma}{original lemmas}
\Crefname{originallemma}{Original lemma}{Original lemmas}
\makeatletter
\newlabel{theorem-openness-flatness}{{theorem-openness-flatness}{source}{}{sourcepointer-4}{}}
\newlabel{theorem-openness-flatness@cref}{{[originallemma][0][]theorem-openness-flatness}{[1][0][]source}{}{}{}}
\makeatother
\title{The locus of points where a finitely presented flat module is flat over an intermediate scheme is open and commutes with base change}
\author{The Stacks Project Authors}
\date{stacks-project:more-morphisms.tex:local-snapshot-20261011; GNU FDL 1.2 or later}
\begin{document}
\maketitle
\noindent\textit{Editorial note.} Statement span is lines 4701-4720 of sources/stacks-more-morphisms/pinned-original-source.tex; proof span is lines 4722-4756, the complete author proof (33 body lines).\par\medskip
\noindent\textit{Editorial note.} Navigation stubs: 4 cited-but-not-included labels are supplied as non-load-bearing pointers with their pinned tags where the snapshot tag index has them.\par\medskip
\noindent\textit{Editorial note.} extra\_packages is empty on purpose: the retained preamble already loads xy with options and a second declaration clashes in the assembler wrapper.\par\medskip
\noindent\textit{Editorial note.} Text is retained verbatim, including the author's own typos and placeholders.\par\medskip
\noindent\textit{Allowed standard prerequisites (editorial).}\begin{itemize}
\item Explicit prerequisite boundary. Flatness of a module over a morphism, finite presentation of modules and morphisms, and the open nature of the flat locus; the base-change comparison is retained.
\end{itemize}\medskip
\ifdefined\hypertarget\hypertarget{sourcepointer-1}{}\fi
\noindent\textit{Source reference.} Original lemma lemma-flat-locus-base-change: sources/stacks-more-morphisms/pinned-original-source.tex, source lines 4468--4486. Cited by the retained statement or proof but not part of the collected spans; supplied only as a navigation pointer. The argument does not depend on its content.\par\medskip
\ifdefined\hypertarget\hypertarget{sourcepointer-2}{}\fi
\noindent\textit{Source reference.} Original lemma morphisms-lemma-finite-presentation-permanence: sources/stacks-more-morphisms/pinned-original-source.tex, source lines 4756--4756. Cited by the retained statement or proof but not part of the collected spans; supplied only as a navigation pointer. The argument does not depend on its content.\par\medskip
\ifdefined\hypertarget\hypertarget{sourcepointer-3}{}\fi
\noindent\textit{Source reference.} Original lemma theorem-criterion-flatness-fibre: sources/stacks-more-morphisms/pinned-original-source.tex, source lines 4590--4599. Cited by the retained statement or proof but not part of the collected spans; supplied only as a navigation pointer. The argument does not depend on its content.\par\medskip
\ifdefined\hypertarget\hypertarget{sourcepointer-4}{}\fi
\noindent\textit{Source reference.} Original lemma theorem-openness-flatness: sources/stacks-more-morphisms/pinned-original-source.tex, source lines 4446--4449. Cited by the retained statement or proof but not part of the collected spans; supplied only as a navigation pointer. The argument does not depend on its content.\par\medskip
\begin{lemma}
\label{lemma-base-change-criterion-flatness-fibre}
Let $S$ be a scheme. Let $f : X \to Y$ be a morphism of schemes over $S$.
Let $\mathcal{F}$ be a quasi-coherent $\mathcal{O}_X$-module.
Assume
\begin{enumerate}
\item $X$ is locally of finite presentation over $S$,
\item $\mathcal{F}$ an $\mathcal{O}_X$-module of finite presentation,
\item $\mathcal{F}$ is flat over $S$, and
\item $Y$ is locally of finite type over $S$.
\end{enumerate}
Then the set
$$
U = \{x \in X \mid \mathcal{F} \text{ flat at }x \text{ over }Y\}.
$$
is open in $X$ and its formation commutes with arbitrary base change:
If $S' \to S$ is a morphism of schemes, and $U'$ is the set of points
of $X' = X \times_S S'$ where $\mathcal{F}' = \mathcal{F} \times_S S'$
is flat over $Y' = Y \times_S S'$, then $U' = U \times_S S'$.
\end{lemma}

\begin{proof}
By
Morphisms, Lemma \ref{morphisms-lemma-finite-presentation-permanence}
the morphism $f$ is locally of finite presentation.
Hence $U$ is open by
Theorem \ref{theorem-openness-flatness}.
Because we have assumed that $\mathcal{F}$ is flat over $S$ we see that
Theorem \ref{theorem-criterion-flatness-fibre}
implies
$$
U = \{x \in X \mid \mathcal{F}_s \text{ flat at }x \text{ over }Y_s\}.
$$
where $s$ always denotes the image of $x$ in $S$. (This description also
works trivially when $\mathcal{F}_x = 0$.) Moreover, the assumptions
of the lemma remain in force for the morphism $f' : X' \to Y'$
and the sheaf $\mathcal{F}'$. Hence $U'$ has a similar description.
In other words, it suffices to prove that given
$s' \in S'$ mapping to $s \in S$ we have
$$
\{x' \in X'_{s'} \mid
\mathcal{F}'_{s'} \text{ flat at }x' \text{ over }Y'_{s'}\}
$$
is the inverse image of the corresponding locus in $X_s$.
This is true by
Lemma \ref{lemma-flat-locus-base-change}
because in the cartesian diagram
$$
\xymatrix{
X'_{s'} \ar[d] \ar[r] & X_s \ar[d] \\
Y'_{s'} \ar[r] & Y_s
}
$$
the horizontal morphisms are flat as they are base changes by the flat
morphism $\Spec(\kappa(s')) \to \Spec(\kappa(s))$.
\end{proof}

\end{document}
